% English translation of checked_pages/esq_119_032.tex; source: cote 119, site 32, PDF page 33, typescript page 30.
% AI-translated by GPT-6 Astra (OpenAI, ChatGPT Pro mode).
\scanpage{119/32}
\mspage{30}

such as the operations of contraction and gluing (without even proceeding
to operations of the type $X \longrightarrow \mathrm{Aut}(X)$, which take one out of the
paradise of finite-dimensional ``spaces''). That is a way of
beating around the bush ! This situation, as so many times already in
the history of our science, simply brings to light the almost
insurmountable inertia of the mind, burdened by conditioning
of considerable weight, when it comes to looking at a
foundational question, and therefore at the very context in which one lives,
breathes and works - instead of accepting it as an immutable
given. It is surely because of this inertia that it took
millennia for an idea or a reality as childlike as
zero, a group or a topological shape to gain citizenship
in mathematics. It is also surely because of it that the straitjacket
of general topology continues to be patiently dragged along by
generations of topologists, ``wildness'' being borne as
an inescapable fate supposedly rooted in the very nature
of things.

My approach toward possible foundations for a tame
topology was an axiomatic one. Rather than declaring
(which would certainly be perfectly reasonable) that
the ``tame spaces'' sought are precisely, say,
Hironaka's semianalytic spaces\frapparatus{V-FR-032-001}, and then developing in
this context the arsenal of constructions and notions familiar in
topology, together, certainly, with those that until now had been
unable to be developed, and for good reason, I preferred to try to bring out
which of the geometric properties of the notion of a semianalytic
set in a space $\mathbb{R}^{n}$ allow such sets to be used as
local ``models'' for a notion of ``\uline{tame space}'' (in this case,
semianalytic), and what (one hopes !) makes this notion of a
tame space sufficiently flexible actually to serve
as a basic notion for a ``tame topology'' capable of expressing
the topological intuition of shapes with ease. Thus, once
the necessary foundational work has been accomplished, there will appear not
\uline{one} ``tame theory'', but a vast infinitude, ranging from the most

\margin{32}

